\documentclass[11pt,letterpaper]{article}
\usepackage[T1]{fontenc}
\usepackage{lmodern,microtype}
\usepackage[margin=0.8in]{geometry}
\usepackage{amsmath,amssymb,amsthm,mathtools,booktabs}
\usepackage[colorlinks=true,linkcolor=blue,urlcolor=blue,citecolor=blue]{hyperref}
\setlength{\parindent}{0pt}
\setlength{\parskip}{3pt}
\AtBeginDocument{%
\setlength{\abovedisplayskip}{8pt}%
\setlength{\belowdisplayskip}{8pt}}
\hypersetup{pdftitle={Two counterexamples to a force-energy bound},
pdfsubject={Exact counterexamples to Herman, Knorn and Ahlen, Theorem 4(iii)}}
\newcommand{\R}{\mathbb R}
\newcommand{\norm}[1]{\left\lVert#1\right\rVert}
\newcommand{\sig}{\sigma_{\min}}
\newtheorem{proposition}{Proposition}
\newtheorem{lemma}{Lemma}
\begin{document}
\begin{center}
{\LARGE Two counterexamples to a force--energy bound\par}
\vspace{5pt}
16 October 2026
\end{center}

The conclusion contradicted here is \textbf{Theorem 4(iii), equation (7)},
of Herman, Knorn and Ahl\'en~\cite{HKA}, printed page 15. We give an analytic
five-follower counterexample and an exact-arithmetic, computer-assisted
counterexample with 128 followers and zero initial state. The latter is an
exact polynomial replacement of the numerical witness, with a certified
relative violation that rounds to $7.26\%$.

\section{The claim and its hypotheses}

Let $U$ have ones on its first superdiagonal and zero elsewhere, and put
\[
 B=I-U,\qquad K=BB^T,\qquad
 D_h=(B+h|B|)B^T.
\]
Here $|B|$ means entrywise absolute value. We take all masses, stiffness
coefficients and damping coefficients equal to one: in the paper's notation
$M=A=R=I$. Position coupling is symmetric, $h_\Delta=0$, and velocity
asymmetry is $h_p=h>0$. With a stationary leader, momentum $p$ and spacing
error $\Delta$, the paper's system (4)--(5) becomes
\begin{equation}\label{eq:system}
 \dot p=-D_hp+B\Delta+d,\qquad \dot\Delta=-B^Tp,
 \qquad H(p,\Delta)=\tfrac12(\norm p^2+\norm\Delta^2).
\end{equation}
All unqualified vector and operator norms are Euclidean. In equilibrium
displacement coordinates, $p=\dot x$ and $\Delta=-B^Tx$, so equivalently
\begin{equation}\label{eq:position}
 \ddot x+D_h\dot x+Kx=d,\qquad
 H=\tfrac12(\norm{\dot x}^2+\norm{B^Tx}^2).
\end{equation}

In this specialization, the claimed bound is
\begin{equation}\label{eq:claim}
 H(t)\leq H(0)+\frac{J}{2\sig(D_h)},\qquad
 J=\int_0^\infty\norm{d(t)}^2\,dt.
\end{equation}
The integral is over the whole half-line. Both forces below vanish after their
stated terminal time $T$, so $J=\int_0^T\norm d^2$.

\begin{center}
\begin{tabular}{@{}p{0.37\linewidth}p{0.57\linewidth}@{}}
\toprule
Hypothesis or convention & Check for both examples\\
\midrule
Assumption 2: uniformly positive, bounded $m_i,a_i,r_i$
 & All are one, for every index and every flock size.\\[3pt]
Assumption 3: $r_i\geq r_{i+1}$
 & Equality holds.\\[3pt]
Symmetric position coupling; $h_p>0$
 & $h_\Delta=0$; $h_p=1$ in \S2 and $h_p=1/2$ in \S3.\\[3pt]
Leader and terminal boundary
 & Leader velocity is zero; the last follower uses only its predecessor,
 exactly as encoded by $B$ in the source.\\[3pt]
Admissible disturbances and initial states
 & Forces have compact support and lie in $L^2$.
 Theorem 4(iii) permits arbitrary initial data; \S3 starts at zero.\\
\bottomrule
\end{tabular}
\end{center}

The problem in the source proof is its equation (8): for nonsymmetric $D_h$,
the relevant quadratic-form lower bound uses
$S=(D_h+D_h^T)/2$, not $\sig(D_h)$. Indeed,
\[
 \dot H=-p^TD_hp+p^Td=-p^TSp+p^Td.
\]
Failure of that proof step alone would not disprove \eqref{eq:claim}. The two examples
below contradict the integrated inequality itself.

\newpage
\section{An analytic counterexample with five followers}

Take $N=5$, $h=1$, so $D_1=2B^T$. Set
\begin{equation}\label{eq:small-data}
 p_*=(3,5,6,5,3)^T,\quad \Delta(0)=0,\quad p(0)=p_*,
 \quad d(t)=\begin{cases}p_*/2,&0\leq t\leq T,\\0,&t>T,\end{cases}
 \quad T=\frac1{20000}.
\end{equation}
In particular, no velocity coupling has a negative weight; the rear velocity
weight is zero when $h=1$.

\begin{lemma}\label{lem:sigma5}
For this matrix, $\sig(D_1)>5/9$, hence $1/(2\sig(D_1))<9/10$.
\end{lemma}
\begin{proof}
Let $C=B^{-1}$, the upper triangular matrix of ones. Then
$(C^TC)_{ij}=\min(i,j)$ and
\[
 \operatorname{tr}\big((C^TC)^2\big)
 =9+28+45+48+25=155.
\]
Since $D_1^{-1}=C^T/2$, the largest eigenvalue of the positive matrix
$D_1^{-1}D_1^{-T}$ satisfies
\[
 \sig(D_1)^{-4}
 \leq \operatorname{tr}\big((D_1^{-1}D_1^{-T})^2\big)=\frac{155}{16}.
\]
Thus $\sig(D_1)^4\geq16/155>(5/9)^4$; the last comparison is the integer
inequality $104976>96875$.
\end{proof}

Direct integer calculations give
\[
 \norm{p_*}^2=104,\qquad p_*^TD_1p_*=28,\qquad
 \norm{p_*/2}^2=26.
\]
Consequently $H(0)=52$, $\dot H(0)=24$, and $J=26T=13/10000$.
Define, only on $[0,T]$,
\[
 F(t)=H(t)-52-\frac9{10}\,26t.
\]
Then $F(0)=0$ and $F'(0)=24-117/5=3/5>0$.
Continuity already proves a violation at some positive time. The following
bound supplies the explicit time in \eqref{eq:small-data}.

Write $y=(p,\Delta)$ and
$\mathcal A=\begin{psmallmatrix}-D_1&B\\-B^T&0\end{psmallmatrix}$.
Since $\norm B\leq2$, $\norm{D_1}\leq4$, we have $\norm{\mathcal A}\leq6$.
Also $\norm{y(0)}<11$ and $\norm{(p_*/2,0)}<6$.
Variation of constants, and $e^u\leq(1-u)^{-1}$ for $0\leq u<1$, yield
\[
 \norm{y(t)}\leq e^{6t}(11+6t)
 \leq\frac{553}{47}<12\qquad(0\leq t\leq T<1/100).
\]
For $0<t<T$, the force is constant and these estimates give
\[
 \norm{\dot y}\leq78,\qquad \norm{\ddot y}\leq468,
 \qquad |\ddot H|\leq78^2+12\cdot468=11700.
\]
Apply these estimates on $[0,T]$. Taylor's integral remainder gives
\begin{equation}\label{eq:small-gap}
 F(T)\geq\frac35T-5850T^2
 =\frac{123}{8000000}>0.
\end{equation}
By Lemma~\ref{lem:sigma5}, the actual claimed right-hand side is strictly
smaller than $52+(9/10)J$. Therefore
\[
 H(T)>52+\frac9{10}J
 >52+\frac{J}{2\sig(D_1)}.
\]
This is a fully analytic counterexample to Theorem 4(iii). The estimates do not
depend on numerical integration or floating-point eigenvalues.

\newpage
\section{An exact counterexample with zero initial state}

\subsection*{How the numerical witness was found}
For $N=128$, $h=1/2$, $T=512$, let
\[
 \mathcal A_x=\begin{pmatrix}0&I\\-K&-D_{1/2}\end{pmatrix},\quad
 F_x=\binom0I,\quad C_x=\operatorname{diag}(B^T,I),\quad
 W_T=\int_0^T e^{s\mathcal A_x}F_xF_x^Te^{s\mathcal A_x^T}\,ds.
\]
The largest possible final energy for a unit-$L^2$ force and zero initial state
is $\tfrac12\lambda_{\max}(C_xW_TC_x^T)$; its top eigenvector determines a
maximizing force.
Numerically this gives energy $43.74150560$ against a claimed coefficient
$40.78143255$. These calculations locate a witness; they are not needed to
verify the exact counterexample below.

\subsection*{An exactly specified trajectory and force}
Put $m=2048$, $t_j=j/4$, $M_0=2^{40}$ and $Q=4M_0$.
The supplied integer data consist of
$X_j,V_j\in\mathbb Z^{128}$ for $0\leq j\leq m$, with
$X_0=V_0=0$. On $[t_j,t_{j+1}]$, set $u=4(t-t_j)$ and
\begin{equation}\label{eq:hermite}
 x(t)=Q^{-1}\sum_{k=0}^3 C_{j,k}u^k,
\end{equation}
where
\begin{align*}
 C_{j,0}&=4X_j,&C_{j,1}&=V_j,\\
 C_{j,2}&=12(X_{j+1}-X_j)-2V_j-V_{j+1},&
 C_{j,3}&=8(X_j-X_{j+1})+V_j+V_{j+1}.
\end{align*}
These formulas give $x(t_j)=X_j/M_0$ and $\dot x(t_j)=V_j/M_0$ on both sides
of every knot. Thus $x$ is $C^1$ and starts from $x(0)=\dot x(0)=0$.
\emph{Define}, almost everywhere on $[0,T]$,
\begin{equation}\label{eq:exact-force}
 d=\ddot x+D_{1/2}\dot x+Kx.
\end{equation}
After $T$, set $d=0$ and continue the homogeneous ODE from the endpoint state.
This force is a bounded piecewise polynomial, hence belongs to $L^2$.
Jumps at the finitely many knots cause no difficulty: $x$ and $\dot x$ are
continuous, so there are no impulsive acceleration terms.

Rounding a numerical trajectory proposed these integer coefficients. They now
define a new trajectory exactly, so no interpolation error enters the proof.

\subsection*{Exact input cost and terminal energy}
Let $D_2=2D_{1/2}$. Its diagonal is $(4,\ldots,4,3)$, its subdiagonal is $-3$,
and its superdiagonal is $-1$. The matrix $K$ has diagonal $(2,\ldots,2,1)$
and adjacent off-diagonals $-1$. With $C_{j,k}=0$ for $k>3$, put
\[
 G_{j,k}=2KC_{j,k}+4(k+1)D_2C_{j,k+1}
             +32(k+1)(k+2)C_{j,k+2}.
\]
Then $d(t)=(2Q)^{-1}\sum_{k=0}^3G_{j,k}u^k$. Consequently
\begin{align}
 J&=\frac1{16Q^2}\sum_{j=0}^{m-1}\sum_{k,l=0}^3
            \frac{G_{j,k}^TG_{j,l}}{k+l+1},\label{eq:J-exact}\\
 H(T)&=\frac{\norm{V_m}^2+\norm{B^TX_m}^2}{2M_0^2}.\label{eq:H-exact}
\end{align}
Integer arithmetic and rational summation give
\begin{align*}
 J&=\frac{16247963052961153388411898917}
          {16247963015620616108051005440},\\
 H(T)&=\frac{6610029438766194687931881}
            {151115727451828646838272}.
\end{align*}
In particular, the exact comparisons
\begin{equation}\label{eq:simpleHJ}
 0<J<\frac{101}{100},\qquad H(T)>43
\end{equation}
already suffice for a contradiction.

\subsection*{A certified singular-value comparison}
For a positive rational $s=a/b$, form the integer pentadiagonal matrix
\[
 \mathcal G_s=b^2D_2^TD_2-4a^2I
             =4b^2\big(D_{1/2}^TD_{1/2}-s^2I\big).
\]
The verifier performs exact rational $LDL^T$ elimination and reconstructs the
factorization. For $s=3/250$, all 128 pivots are strictly positive. Hence
$\mathcal G_s\succ0$ and $\sig(D_{1/2})>3/250$. Combining this with
\eqref{eq:simpleHJ},
\begin{equation}\label{eq:zero-contradiction}
 \frac{J}{2\sig(D_{1/2})}
 <\frac{125}{3}J
 <\frac{505}{12}
 <43<H(T).
\end{equation}
Since $H(0)=0$, this explicitly contradicts \eqref{eq:claim}.
If an exactly unit-$L^2$ input is preferred, divide $x,\dot x,d$ by $\sqrt J$;
the inequality is homogeneous, so the contradiction is preserved.

For completeness, the banded elimination uses only the following recurrence.
Writing $g_{ij}=(\mathcal G_s)_{ij}$, $\ell_{ii}=1$, and omitting terms with
negative indices,
\begin{align*}
 \ell_{i,i-2}&=g_{i,i-2}/\delta_{i-2},\\
 \ell_{i,i-1}&=(g_{i,i-1}-\ell_{i,i-2}\delta_{i-2}\ell_{i-1,i-2})/\delta_{i-1},\\
 \delta_i&=g_{ii}-\ell_{i,i-1}^2\delta_{i-1}
                     -\ell_{i,i-2}^2\delta_{i-2}.
\end{align*}
All divisions and signs are evaluated as rational numbers. The certificate
lists every pivot, and the checker verifies the identity
$\mathcal G_s=L\operatorname{diag}(\delta_i)L^T$.

\subsection*{Certifying the size of the violation}
The same exact elimination gives the sharper bracket
\[
 \frac{245209}{20000000}<\sig(D_{1/2})<\frac{24521}{2000000}.
\]
The lower matrix has 128 positive pivots. The upper matrix has 127 positive
and one negative pivot, so Sylvester's law of inertia proves the upper bound.
Using the exact $H(T)$ and $J$ above, these brackets certify
\[
 7.2581\%
 <\frac{H(T)}{J/(2\sig(D_{1/2}))}-1
 <7.25855\%.
\]
Thus $7.26\%$ is a justified rounding for the exact polynomial witness,
independently of floating-point singular-value computations.

\section{Verification, data and scope}

The accompanying certificate bundle contains the complete integer trajectory
and two verifiers. From the repository root, using only the Python standard
library, run
\begin{verbatim}
python3 certificates/energy-short-time/verify.py
python3 certificates/energy-zero-state/verify.py
\end{verbatim}
The first checks the finite integer identities and rational norm/Taylor bounds
in \S2. The second checks the data, zero initial conditions, knot matching,
exact polynomial integrals, every factorization pivot, and the strict energy
comparisons in \S3. Neither requires NumPy, SciPy, an ODE solver, a matrix
exponential or a floating-point eigensolver. Decimal summaries are not used
in any proof decision. An optional proposal script explains the numerical
origin of the data and is never imported by either checker.

The integer data format is one ASCII JSON header followed by 2049 rows of decimal
integers, in the order
\[
 X_{j,1},\ldots,X_{j,128},V_{j,1},\ldots,V_{j,128}.
\]
The compressed file is \texttt{trajectory.txt.gz}; its SHA-256 is
\begin{center}\footnotesize\ttfamily
b2411735af07bf15f213da435844786c417ae3d180760e8948421b304fe2c97f.
\end{center}
The hash identifies the finite witness; correctness follows from the exact
checks, not from a claim that the proposed data are accurate samples.
The five-follower argument is entirely analytic. The zero-state argument is
a computer-assisted proof by exact rational arithmetic, with its data and
checking algorithm supplied. It is not a formalization in a proof assistant.

Smooth compactly supported disturbances also exist: the terminal state depends
continuously on $d$ in $L^2(0,T)$ through the bounded map
$d\mapsto\int_0^T e^{(T-s)\mathcal A_x}F_xd(s)\,ds$.
Thus the strict violation persists under sufficiently small smooth $L^2$
approximations. The explicit certificates use the polynomial forces above.

These examples refute the stated finite coefficient in Theorem 4(iii).
They do not refute the separate asymptotic stability and passivity assertions
in parts (i)--(ii), or the singular-value estimates of Lemma 6.
For a positive definite symmetric part $S=(D_h+D_h^T)/2$, a valid bound follows
directly from the energy identity and Young's inequality:
\[
 H(t)\leq H(0)+\frac1{4\lambda_{\min}(S)}
                     \int_0^t\norm{d(s)}^2\,ds.
\]
Indeed, $-\lambda\norm p^2+p^Td\leq\norm d^2/(4\lambda)$.
This supplies a valid replacement estimate; it does not establish the
paper's asserted scaling in $N$ for the sharper, system-dependent gain.

\begin{thebibliography}{1}
\bibitem{HKA}
I.~Herman, S.~Knorn and A.~Ahl\'en,
\emph{Disturbance scaling in bidirectional vehicle platoons with different
asymmetry in position and velocity coupling},
Automatica \textbf{82} (2017), 13--20.
\href{https://doi.org/10.1016/j.automatica.2017.04.010}
{doi:10.1016/j.automatica.2017.04.010}.
Theorem 4 and equations (3)--(9) are on printed p.~15 in the supplied
published PDF; the corresponding author manuscript is dated 15 March 2017.
\end{thebibliography}
\end{document}
